\documentclass{article}
\usepackage[utf8]{inputenc}
\usepackage{amsmath}
\usepackage{geometry}
\geometry{margin=2cm}
\begin{document}

\section*{Questão 118 — ENEM 2024 — Ciências da Natureza}

\subsection*{Interpretação e Estratégia}

A questão apresenta um sistema óptico formado por dois espelhos côncavos, chamados $E_1$ e $E_2$, arranjados de tal forma que o vértice de um coincide com o foco do outro e vice-versa. Um objeto está posicionado sobre o vértice de $E_1$ e, através da reflexão em ambos os espelhos, forma-se uma imagem tridimensional visível por um observador. O objetivo é determinar a natureza dessa imagem — se é real ou virtual — e a distância vertical entre um ponto do objeto e seu ponto imagem correspondente.

A estratégia para resolver consiste em:
\begin{itemize}
  \item Reconhecer as características dos espelhos côncavos e suas relações entre vértices e focos.
  \item Analisar o percurso dos raios (que refletem em $E_2$ e depois em $E_1$) para compreender o local de formação da imagem.
  \item Usar as dimensões apresentadas — abertura de 5 cm, distância de 3,8 cm para cada espelho — para calcular a distância vertical entre objeto e imagem.
  \item Compreender que a natureza real da imagem decorre da convergência física dos raios neste sistema.
\end{itemize}

\subsection*{Conteúdo e Mini revisão}

\begin{itemize}
  \item \textbf{Espelhos côncavos:} podem formar imagens reais ou virtuais, dependendo da posição do objeto em relação ao foco e ao centro de curvatura.
  \item \textbf{Imagem real vs. virtual:} imagem real é formada quando os raios convergem realmente; virtual, quando os raios parecem divergir de um ponto.
  \item \textbf{Interpretação de medidas:} a soma das distâncias verticais relevantes no esquema é crucial para determinar o deslocamento entre objeto e imagem.
\end{itemize}

\subsubsection*{Resumo teórico}

Espelhos côncavos formam imagens reais quando o objeto está fora do foco, permitindo que os raios refletidos se encontrem em um ponto físico (imagem). Imagens virtuais são aquelas que o observador percebe como localizadas em um ponto de onde os raios parecem divergir. O sistema apresenta um arranjo especial dos espelhos, cujo efeito é a criação de uma imagem tridimensional real.

\subsection*{Resolução detalhada e Pulo do Gato}

\begin{enumerate}
  \item Observamos que os espelhos $E_1$ e $E_2$ são côncavos e estão dispostos de modo que o vértice de $E_1$ coincide com o foco de $E_2$, e vice-versa.
  \item O objeto está colocado sobre o vértice de $E_1$.
  \item Os raios de luz saem do objeto, refletem primeiro em $E_2$ e depois em $E_1$, convergindo em pontos que formam a imagem tridimensional.
  \item A imagem, portanto, é real, pois os raios convergem fisicamente.
  \item As dimensões indicam uma distância vertical de $3,8\,\mathrm{cm}$ de cada espelho até o centro, então a distância total vertical entre objeto e imagem é

  \[
  3{,}8\,\mathrm{cm} + 3{,}8\,\mathrm{cm} = 7{,}6\,\mathrm{cm}.
  \]

  \item Concluímos que a distância vertical entre cada ponto do objeto e seu correspondente ponto imagem é $7,6\,\mathrm{cm}$, e que a imagem é real.
\end{enumerate}

\textbf{Pulo do gato:} a soma das distâncias verticais de cada espelho indica a distância total entre o objeto e a imagem, e a convergência real dos raios de luz garante que a imagem é real e não virtual.

\subsection*{Armadilhas e Alternativas}

\begin{tabular}{|c|p{4cm}|p{4cm}|p{5cm}|p{5cm}|}
  \hline
  \textbf{Alternativa} & \textbf{Raciocínio provável} & \textbf{Tipo de erro} & \textbf{Por que é enganoso} & \textbf{Por que é incorreto} \\
  \hline
  A & Acreditar que a distância vertical é a abertura de 5 cm & Confusão entre distância e abertura & A abertura aparece no esquema em destaque & A distância vertical real é 7,6 cm, soma das duas partes de 3,8 cm \\
  \hline
  B & Usar apenas um trecho de 3,8 cm como distância total & Uso parcial das medidas & 3,8 cm é destacado no esquema & Precisa somar os dois trechos para chegar a 7,6 cm \\
  \hline
  D & Pensar que imagem é virtual por ser tridimensional & Confusão de percepção e natureza da imagem & Imagens 3D são frequentemente associadas a imagens virtuais & Raios convergem fisicamente, logo imagem é real \\
  \hline
  E & Combinar distância incorreta e natureza virtual & Erro duplo na análise & Combina valores conhecidos com conclusões erradas & A imagem correta é real e a distância correta é 7,6 cm \\
  \hline
\end{tabular}

\subsection*{Código da Questão}

CN_OPTICA_3D_001

\end{document}